%%% Introduction --- 
%% 
%% Filename: Intro-cs.tex
%% Author: Jacob Ringfjord & Max Wilén
%%% Code:


\chapter{Introduction}
\label{cha:introduction}

\begin{comment}
SULEMAN COMMENT:
    "Good job.
    Quick feedback!
    1. Add couple of more citions plus see the ACM survey paper I shared with you in very first meeting so you will get the idea how to start the first para of intro.."
\end{comment}

The rise of new technology and the growing number of connected devices has led to an even greater variety of software applications. Adding to this, the wide range of available development tools and software products enables developers to release projects more quickly, leading to faster development cycles and shorter intervals between software releases \cite{ACM_Current_and_Future_Trends_in_Mobile_Device_Forensics_A_Survey}. While this approach enhances security and reliability for companies, it also presents challenges in forensic analysis. Investigations into these applications take a longer time, and the sheer volume of applications further compounds the problem \cite{android_matching, introductionForensics}. Moreover, the rapid growth of artificial intelligence in recent years has led to an increase in the number of cyberattacks \cite{gurtov_3}, increasing the workload for forensics, making it essential for tools used during investigations to evolve accordingly.

This thesis further explores this problem and presents a novel method for automating version comparison and categorizing relevant changes for use in reverse engineering during forensic investigations.

\section{Motivation}
\label{sec:motivation}

\begin{comment}
SULEMAN COMMENT:
    "Motivation means why in research for example why you are motivated to work on this issue or why you to propsoe solution for this. See some gaps in exisitng literature get motivated and why exisitng solution can not fix these valuerabilties. Last para should tell okay you are doing this and in will do this and it will improve these  issues. Further i will tell you in the meeting"

SULEMAN COMMENT:
    "Nice work! However, i will say here you have only one paper instead you can cite 3 or 4 papers here if possible and get motivation from these papers like what issues or challeneges we have and these challeneges have that impact so therefore you are doing this this work in this thesis"
\end{comment}

Applications of digital forensic interest are often analyzed to examine functionality that may affect the results of digital forensic investigations. Therefore, the forensic process relies on prior analysis of the application to have an accurate understanding of the software, ensuring that evidence is correctly extracted, interpreted, and admissible in court. The method chosen in a forensic investigation is dependent on what type of media and application the digital evidence is to be extracted from. When a new version of an application is released, valuable knowledge about its functionality may no longer hold for reverse engineers \cite{android_matching}. Even minor software updates may change how the information is gathered and which forensic methods should be applied \cite{acm_fingerprinting_survey, nij2004forensic}. As a result, the newly released software version may require manual analysis, for instance, performing reverse engineering.

Given the rapidly evolving industry in software development and the increasing number of available applications, the pace needed to keep up with the release cycles becomes challenging. For instance, the open-source end-to-end encrypted messaging application Signal is a frequent example of an application of forensic interest. It has multiple software releases every week \cite{githubSignal}. Other encrypted messaging applications monitored by authorities are EncroChat, Exclu, and Ghost \cite{europaInternationalOperation}, which also have frequent release cycles but are not open-source.

However, this evolution is not limited to messaging platforms. As the number of connected devices increases and their computational capabilities grow, the technologies behind secure communications are expected to extend into other domains, such as vehicles \cite{gurtov_1}, infrastructure systems, and healthcare \cite{gurtov_2}. \textcite{gurtov_1} highlights that the adoption of emerging secure communication technologies in the automotive sector brings a new wave of security and privacy challenges. This shift broadens the landscape of applications with forensic interest as the focus excels from encrypted messaging applications to entire ecosystems of interconnected devices and systems.

For digital forensics, an unreasonable amount of work is needed to conduct a full analysis of every version of all relevant software on these devices. As described by the thesis client, without knowledge of the relevant changes between two software releases, forensic work becomes an ad hoc process in which all differences must be investigated, rather than focusing on those relevant to the forensic investigation. By categorizing changes and filtering them based on their relevance to the forensic investigation, a decrease in the time needed to analyze a single release of an application could be possible.

%%%%%%% KÄNS LITE FÖR MYCKET AIM
% This thesis investigates the development of a method to automate the digital forensic process by systematically analyzing differences between software versions.

%%%%%%%% TEXT TAGEN FRÅN DIGITAL FORENSIC METHODS
% As technology continues to advance and digital media becomes embedded in all aspects of modern society, the challenge of managing the growing volume of digital evidence requiring forensic analysis to ensure its admissibility in court has become increasingly significant \cite{introductionForensics}.


% This thesis aims to develop a method for comparing the source code between different versions of the same application. The comparison is intended to help assess whether the changes in the code are significant enough to warrant a manual analysis. This could potentially be achieved through symbolic execution or semantic analysis.

\begin{comment}
\section{Motivation}
\label{sec:motivation}

This is where the studied problem is described from a general
point of view and put in a context which makes it clear that
it is interesting and well worth studying. The aim is to make
the reader interested in the work and create an urge to
continue reading.

Basic points to cover:
\begin{itemize}
\item What is the background for this problem?
\item Why is it interesting/important to study this?
\item What are the challenges? Why is it hard?
\item What does the existing literature say about this problem? What is already known?
\item What are the knowledge gaps?
\end{itemize}
\end{comment}

\section{Aim}
\label{sec:aim}

This thesis aims to develop a novel method for automating version comparison and categorization of relevant changes for use in tandem with reverse engineering of software applications in forensic analysis.

To achieve this, the thesis investigates the development of a tool that can detect programmatic changes between two versions of a software application, categorize these changes, and later filter them based on relevance in the forensic investigation. All without having knowledge of the source code, and only analyzing the application binaries.

\section{Research questions}
\label{sec:research-questions}

A set of research questions has been defined to answer questions forensic investigators may ask when examining software with different versions.

\begin{enumerate}

    \item What current approaches exist to automatically identify version changes in Android application binaries to support forensic analysis?
    
    \item How to effectively categorize and rank detected changes without access to source code?

    \item What is the effectiveness of regular expression-based categorization in filtering forensic-relevant changes in binaries?

\end{enumerate}



\section{Approach}
\label{sec:approach}

\begin{comment}
SULEMAN COMMMENT:
    "only highlight those things which will be covered in this thesis dont put light on those stuff or write which is not part of this thesis"
\end{comment}

This thesis addresses the research questions outlined in Section~\ref{sec:research-questions} by exploring binary analysis techniques and the tools commonly used for analyzing software executables. Following this, a deep dive into categorizing the discovered changes based on functionality defined as relevant, and the method to perform this is considered the most effective. Ultimately, the research seeks to determine how changes in software applications may impact the forensic analysis process.

This thesis's primary focus is on producing a method for assessing whether the changes between the two versions' core operations are significant enough to warrant manual reverse engineering of the application.

\section{Limitations}
\label{sec:limitations}
The thesis does not cover forensic methods or tools used in the analysis, assuming they are functional and relevant to the core operations of the examined software. During a forensic analysis, various data sources such as RAM dumps, network traffic, and device storage may be examined, which requires specialized tools or further reverse-engineering to extract information. While certain application changes can influence how such tools are applied, this thesis only focuses on presenting and categorizing software changes to maintain a limited scope. Therefore, the usage of the output from the developed tool is never discussed.

\section{Contributions}
\label{sec:contributions}
\begin{comment}    
- research about forensic methods in general
- how changes in software source code affects the application of these forensic methods
- research on differences between semantic analysis and symbolic execution
    - how semantic analysis and symbolic execution can be applied to effectively determine these software changes.
- create a experimental scenario/method for when one of these (semantic analysis or symbolic execution) can be used, and applied for work in a forensic analysis
- present findings and reason for one of the methods being chosen.
\end{comment}

The presented findings can effectively be used to quickly determine the significance of changes between two different versions of the same software. The tool is not intended to replace manual reverse engineering or forensic methods, but rather to support the process by helping analysts identify what changes are worthy of further analysis. This can indicate whether the existing forensic methods, assumptions, and knowledge remain valid or need to be updated. The main contributions of this thesis are to provide: 

\begin{enumerate}
    \item \textbf{Research:} Present findings in current binary analysis methods and how these are applied in a version comparison detection setting, followed by effectively categorizing these changes.

    \item \textbf{Experiment:} An experimental method, in the form of a developed Python program, that can be applied when detecting and categorizing changes between two software versions, without knowing the source code.
\end{enumerate}


%%%%%%%%%%%%%%%%%%%%%%%%%%%%%%%%%%%%%%%%%%%%%%%%%%%%%%%%%%%%%%%%%%%%%%
%%% Intro.tex ends here